\begin{definition}[4.6.9]
\label{IV.2.4.6.9-fr}
Sean $k$ un cuerpo, $X$ un $k$-preesquema y $x$ un punto de $X$. Se dice que
$X$ es \emph{geométricamente reducido} (o \emph{separable}) \textup{(resp.}
\emph{geométricamente puntualmente íntegro}) en el punto $x$ sobre $k$ si,
para toda extensión $k'$ de $k$ y todo punto $x'$ de
$X'=X\otimes_k k'$ situado sobre $x$, $X'$ es reducido \textup{(resp.} íntegro) en
$x'$ \textup{(es decir,} $\mathcal O_{X',x'}$ es reducido \textup{(resp.} íntegro)). Se dice que
$X$ es \emph{geométricamente puntualmente íntegro} si $X$ es geométricamente
puntualmente íntegro en todos sus puntos, es decir, si para toda extensión
$k'$ de $k$ todos los anillos locales de $X\otimes_k k'$ son íntegros \textup{(en
cuyo caso se dice también que $X\otimes_k k'$ es} puntualmente íntegro\textup{)}.
\end{definition}

Observemos que, para que $X$ sea geométricamente reducido sobre $k$ (4.6.2),
es necesario y suficiente que sea geométricamente reducido sobre $k$ en todo
punto $x\in X$.

\begin{proposition}[4.6.10]
\label{IV.2.4.6.10-fr}
Sean $k$ un cuerpo, $X$ un $k$-preesquema, $k'$ una extensión de $k$, $x'$ un
punto de $X'=X\otimes_k k'$ y $x$ su imagen en $X$. Para que $X$ sea
geométricamente reducido \textup{(resp.} geométricamente puntualmente íntegro) en $x$
sobre $k$, es necesario y suficiente que $X'$ lo sea en $x'$ sobre $k'$. En
particular, para que $X$ sea geométricamente puntualmente íntegro sobre $k$,
es necesario y suficiente que $X'$ sea geométricamente puntualmente íntegro
sobre $k'$.
\end{proposition}

\begin{proof}
Basta demostrar la primera afirmación. Como la condición es evidentemente
necesaria, demostremos que es suficiente. Supongamos, pues, que $X'$ es
geométricamente reducido (resp. geométricamente puntualmente íntegro) en el
punto $x'$ sobre $k'$ y demostremos que $X$ lo es en el punto $x$ sobre $k$,
es decir, que para toda extensión $k''$ de $k$ y todo punto $x''$ de
$X''=X\otimes_k k''$ situado sobre $x$, $\mathcal O_{X'',x''}$ es reducido
(resp. íntegro). Para ello, observemos que existe un punto $z$ de
$X'\times_X X''$ que se proyecta sobre $x'$ y $x''$ (\textbf I, 3.4.7). Sea $s$ el
punto de $\operatorname{Spec}(k'\otimes_k k'')$ imagen de $z$ (\textbf I, 3.4.9);
pongamos $K=\boldsymbol{k}(s)$, $Z=X\otimes_k K$, de modo que $K$ puede considerarse como
extensión compuesta de $k'$ y $k''$, y $z$ como un punto de $Z$ cuya imagen
en $X'$ (resp. $X''$) es $x'$ (resp. $x''$). Como $X'$ es geométricamente
reducido (resp. geométricamente puntualmente íntegro) en el punto $x'$ sobre
$k'$, el anillo $\mathcal O_{Z,z}$ es reducido (resp. íntegro), y como
$\mathcal O_{Z,z}$ es un $\mathcal O_{X'',x''}$-módulo fielmente plano, se
sigue que $\mathcal O_{X'',x''}$ es reducido (resp. íntegro)
$(\mathbf{0}_{\mathrm I},6.5.1)$.
\end{proof}

\begin{corollary}[4.6.11]
\label{IV.2.4.6.11-fr}
Supongamos que $k$ es perfecto \textup{(resp.} algebraicamente cerrado). Para que $X$
sea geométricamente reducido \textup{(resp.} geométricamente puntualmente íntegro) en
el punto $x$ sobre $k$, es necesario y suficiente que $\mathcal O_{X,x}$ sea
reducido \textup{(resp.} íntegro).
\end{corollary}

\begin{proof}
Es evidentemente necesario. Recíprocamente, si se cumple esta condición,
entonces, para toda extensión $k'$ de $k$,
$\operatorname{Spec}(\mathcal O_{X,x})\otimes_k k'$ es reducido (4.6.4)
(resp. íntegro (4.4.4));

\oldpage[IV-2]{71}

por tanto, para todo punto $x'$ de $X'=X\otimes_k k'$ situado sobre $x$,
$\mathcal O_{X',x'}$, que es isomorfo a un anillo local de
$\operatorname{Spec}(\mathcal O_{X,x})\otimes_k k'$, es reducido (resp.
íntegro).
\end{proof}

\begin{proposition}[4.6.12]
\label{IV.2.4.6.12-fr}
Sean $k$ un cuerpo; $X$ un $k$-preesquema, $x$ un punto de $X$ y $\Omega$ una
extensión perfecta de $k$. Las condiciones siguientes son equivalentes:
\begin{enumerate}
\item[\textup{a)}] $X$ es geométricamente reducido sobre $k$ en el punto $x$, es
decir, para toda extensión $k'$ de $k$ y todo punto $x'$ de
$X'=X\otimes_k k'$ situado sobre $x$, $\mathcal O_{X',x'}$ es reducido.
\item[\textup{b)}] El preesquema $X\otimes_k\Omega$ es reducido en un punto situado
sobre $x$.
\item[\textup{c)}] Para toda extensión finita radicial $k'$ de $k$,
$X'=X\otimes_k k'$ es reducido en el único punto situado sobre $x$.
\item[\textup{d)}] $\operatorname{Spec}(\mathcal O_{X,x})$ es geométricamente reducido
sobre $k$.
\item[\textup{e)}] $\mathcal O_{X,x}$ es reducido y, para toda componente irreducible
$Z$ de $X$ que contiene a $x$, de punto genérico $z$, $\boldsymbol{k}(z)$ es una extensión
\emph{separable} de $k$.
\end{enumerate}
\end{proposition}

\begin{proof}
Las implicaciones \emph{d)} $\Rightarrow$ \emph{a)} $\Rightarrow$ \emph{b)} $\Rightarrow$ \emph{c)} son
inmediatas, teniendo en cuenta para la última (4.6.10), (4.6.11) y el hecho de
que toda extensión radicial de $k$ es isomorfa a una subextensión de
$\Omega$. Demostremos que \emph{c)} implica \emph{d)}. Basta mostrar que, para toda
extensión radicial finita $k'$ de $k$, $\mathcal O_{X\,x}\otimes_k k'$ es
reducido (4.6.1); ahora bien, este anillo es el anillo local de
$X'=X\otimes_k k'$ en el único punto $x'$ de $X'$ situado sobre $x$
(\textbf I, 3.5.8 y 3.5.7); es, por tanto, reducido en virtud de la hipótesis \emph{c)}. Por
último, la equivalencia de \emph{d)} y \emph{e)} resulta de la equivalencia de \emph{b)} y \emph{e)} en
(4.6.1), pues las componentes irreducibles de $X$ que contienen a $x$ se
corresponden biyectivamente con las componentes irreducibles de
$\operatorname{Spec}(\mathcal O_{X,x})$ (\textbf I, 2.4.2).
\end{proof}

\begin{corollary}[4.6.13]
\label{IV.2.4.6.13-fr}
Bajo las hipótesis de \textup{(4.6.12)}, supongamos además que $X$ es localmente
noetheriano. Entonces las condiciones a) a e) de \textup{(4.6.12)} son también
equivalentes a la siguiente:
\begin{enumerate}
\item[\textup{f)}] Existe una vecindad abierta $U$ de $x$ que es geométricamente
reducida sobre $k$.
\end{enumerate}
\end{corollary}

\begin{proof}
En efecto, esta condición implica trivialmente que $X$ es geométricamente
reducido en el punto $x$ sobre $k$. Recíprocamente, si así ocurre, como el
anillo $\mathcal O_{X,x}$ es reducido, existe una vecindad abierta $U$ de $x$
en $X$ que es reducida (\textbf I, 6.1.13); tomando $U$ suficientemente pequeña,
puede suponerse además que $U$ no corta ninguna de las componentes
irreducibles de $X$ que no contienen a $x$. Entonces el criterio (4.6.12, \emph{e)})
demuestra que $U$ es geométricamente reducido sobre $k$, teniendo en cuenta
el criterio (4.6.1, \emph{e)}).
\end{proof}

Se sigue de (4.6.13) que, cuando $X$ es localmente noetheriano, el conjunto de
los $x\in X$ en los que $X$ es geométricamente reducido sobre $k$ es
\emph{abierto}, y es el mayor abierto de $X$ geométricamente reducido sobre
$k$.

\begin{env}[4.6.14]
\label{IV.2.4.6.14-fr}
Se sigue de (4.6.10) y (4.6.11) que, si $\Omega$ es una extensión
algebraicamente cerrada de $k$, es equivalente decir que $X$ es
geométricamente puntualmente íntegro en el punto $x$ o que
$X\otimes_k\Omega$ es íntegro en \emph{un} punto situado sobre $x$.

Para que $X$ sea geométricamente puntualmente íntegro en el punto $x$, es
necesario que $X$ sea íntegro en el punto $x$ y que, si $z$ es el punto
genérico de la única componente irreducible de $X$ que contiene a $x$, $\boldsymbol{k}(z)$
sea una extensión separable de $k$; esto se sigue de (4.6.10). Pero estas
condiciones \emph{no son suficientes}, como se desprende del ejemplo (4.5.12,
(ii)). Una condición suficiente, pero no necesaria, es que $X$ sea
geométricamente reducido en el punto $x$ y solo pertenezca a una componente
irreducible de $X$, que

\oldpage[IV-2]{72}

se supone además \emph{geométricamente irreducible}; si, además, $X$ es
localmente noetheriano, existe entonces una vecindad abierta de $x$ en $X$ que
es \emph{geométricamente íntegra}.
\end{env}

Recordemos que, si un preesquema \emph{localmente noetheriano} es puntualmente
íntegro, es \emph{localmente íntegro} (\textbf I, 6.1.13). Si un $k$-preesquema $X$,
localmente de tipo finito sobre $k$, es geométricamente puntualmente íntegro
sobre $k$, se sigue, pues, de esta observación que, para toda extensión $k'$
de $k$, $X\otimes_k k'$ es \emph{localmente íntegro} (en este caso se dice
también que $X$ es \emph{geométricamente localmente íntegro}).

